\section{Ghosting a File}\label{sec:ghost}

\subsection{Accepting the original}
GhostVault accepts only a regular file with a single name. Symbolic links, folders, devices, pipes
and sockets are refused, and so are Windows symbolic links and mount points; a file with more than
one hard link is refused because its data would stay on disk under the other name after
GhostVault removed this one; and nothing inside the data folder can be ghosted. The file is opened
once, without following links (\code{O\_NOFOLLOW}, \code{O\_NONBLOCK}), and the identity of the open
file (device, inode, size, modification time) must match that of the path, so a file replaced at
the moment of opening is refused.

Everything is read through that one handle. The size recorded is the number of bytes actually read,
and after the last chunk the file's size and modification time are compared with those at opening.
If anything changed, nothing is committed and the partial upload is undone. On Windows the handle
is opened with Python's default sharing, which does not include deletion, so while the upload runs
Windows refuses to rename, replace or delete the original; writing into it remains possible and is
detected.

\subsection{The plan}
Before anything is sent, GhostVault draws a random file id and a random key, decides where every
chunk goes, and gives every chunk, and the preview if there is one, a random path of the form
\code{ghostvault/<2 hex>/<32 hex>.bin}. Placement follows one of two strategies, chosen in Settings
or for a single file:
\begin{itemize}
\item \emph{distributed} (the default): the eligible repositories are put in a random order and dealt
chunks in turn, so each of $r$ repositories holds $\lfloor n/r\rfloor$ or $\lceil n/r\rceil$ of a file's
$n$ chunks;
\item \emph{consolidated}: every chunk goes to one eligible repository, chosen at random.
\end{itemize}
Neither strategy is redundancy: every chunk is stored once. A distributed file needs every
repository that holds one of its chunks; a consolidated file needs one repository, which holds all
of it. The record lists the repository of every chunk, and every later operation works from that
list, so changing the setting never moves an existing ghost. With a single eligible repository,
distributed placement cannot distribute; the record and the result message say so.

The plan, the key and the original's identity are then written as the record, in state
\state{uploading}, before the first request (\figref{fig:lifecycle}). From that moment every path the
ghost could ever create on GitHub is known to the journal.

\subsection{Upload and commit}
Blobs can be created only in a repository that has a commit, so an empty repository is first
initialised with an empty file, \code{ghostvault/.keep}. Then the file is streamed: each chunk is
read, padded if it is the last, encrypted, hashed (SHA-256 and its Git blob id), and uploaded as a
blob. Git derives a blob's id from the bytes it holds~\cite{git-objects}, so the id in GitHub's
reply must equal the one computed locally. Only one chunk is held in memory at a time, and nothing is written to disk. A
blob is part of no branch until a commit refers to it, so an interrupted upload leaves only
unreferenced, encrypted objects.

Once every chunk is uploaded, every blob id is written to the record, and only then is anything
committed. Each repository gets one commit: a tree built on the branch's current tree with the new
paths, a commit with the branch head as its parent and the message \code{Update}, and a branch
update that must be a fast-forward. If a commit fails without a clear answer, GhostVault reads the
branch to see what arrived before deciding (\secref{sec:uncertain}).

\subsection{Verification before deletion}\label{sec:verify}
Before the original may be deleted, three things must hold for every repository used:
\begin{enumerate}
\item each blob id was computed locally from the exact ciphertext, and GitHub answered the upload
with the same id;
\item the branch, read again from its head, maps every recorded path to exactly that blob id. Git ids
are content hashes, so this establishes that the branch refers to exactly the uploaded bytes, for
every chunk. A path that is not yet listed is looked for again after 1, 2 and 4\,s, because a read
may be served by a replica that has not caught up; still missing, the ghost fails and is undone;
\item the first chunk and one other chunk chosen at random are downloaded through the same request a
restore uses, and their SHA-256 must equal the uploaded ciphertext's. This shows that the data can be
read back with the current token, not only that it is referenced.
\end{enumerate}
Downloading every chunk again would double the transfer for little more assurance, since the second
check already covers every chunk and GitHub checks object ids itself. What no check made at this
time can rule out is damage or loss on GitHub's side afterwards. The record then moves to
\state{stored}.

\subsection{Removing the original}\label{sec:remove}
Checking a path and then deleting the same path would leave a window in which another program, an
editor saving by rename for example, replaces the file, and the newer file, which was never
uploaded, would be deleted. GhostVault therefore moves the original aside before checking it
(\lstref{lst:remove}). From the rename on, whatever appears under the original's name is never
touched.

\begin{lstlisting}[caption={Removing the original, once the record says \state{removing-original}.},
label={lst:remove}]
file at the path not the uploaded one?    keep it
  (device, inode, size, mtime)
hidden = ".<name>.<12 random hex>.ghostvault-removing"
record hidden in the record (durably)
rename path -> hidden          (atomic, same folder)
open hidden    (Windows: no sharing but reading)
  identity changed?            put back, keep
  SHA-256 != uploaded?         put back, keep
  size or mtime changed now?   put back, keep
  delete     (POSIX: unlink, if hidden still names
              the open file; Windows: by the handle)
put back = publish under the old name, never
  replacing; if taken: "<name> (kept by GhostVault)"
\end{lstlisting}

On Windows the hidden file is opened with read and delete access, sharing only reading, and deleted
through that same handle (\code{SetFileInformationByHandle}): no other program can write to it,
rename it or delete it while it is checked, and opening fails if one already has it open for
writing, so the file deleted is exactly the file checked. Linux and macOS offer no way to delete a
file only if nobody else has it open for writing. There, a program that already had the file open
can still write to it in the moment between the last check and the deletion; GhostVault narrows
that window by checking size and modification time once more just before it, and claims no more.

If any step fails, the original is kept, or put back without replacing anything, and the ghost is
completed as \state{ready} with the reason; Windows answers such as ``access denied'' or ``in use''
are handled the same way. The ghost is complete either way: only the deletion of the original was
given up.
